% !TeX root = main.tex
% Changed: 按电荷源、介质与接收响应重组绪论；保留已核引文，缺项留空。
% Protected: 既有参考文献键、封面题名、实测数据。
% Meaning-Check: NEEDS-LLM
% Risk-Flags: none
\chapter{绪论}
\section{研究背景与意义}
\subsection{水下目标探测的需求与多物理场途径}
随着遥感、海洋测绘和水下装备的发展，观测范围从陆地和空域拓展至远洋和深海。海洋资源调查需要获取水体与海底的信息，海上交通和水下设施维护则需要掌握航行器与作业装备的活动。随着作业范围向深远海延伸，观测任务也从单次调查扩展到对海洋环境和目标活动的持续监测。欧盟委员会发布的海洋政策文件将海洋研究、创新和海上安全列为重点\cite{euOceanPact2025}。我国《“十五五”规划纲要》部署新一轮海洋综合调查和深海工程，将提升深海探测能力列为具体任务\cite{chinaPlan2026}。《新时代的中国国防》白皮书也将近海防御和远海防卫列为海军建设的战略要求\cite{chinaDefense2019}。在这些发展与安全需求下，海洋观测需要掌握海面船舶和水下航行器的位置与活动。

水声探测的发展为水下观测建立了技术基础。20 世纪初的水下声学装置已利用回波开展测深和测距\cite{fessenden1923}。此后，利用主动发射声波获得回波和利用水听器记录目标辐射声，形成了两类水声观测方式。随着传感器及观测平台的发展，声呐可用于水体和海底测绘，也可搭载在自主水下航行器上开展局部观测\cite{noaaAuv}。水声信号的传播受到海水声速分布和海面、海底反射的影响，接收信号包含目标与传播环境的共同作用\cite{gebbie2014}。

水下航行器与周围水体的相互作用还会形成磁场、电场、压力和尾流等特征，水下观测也相应发展出声学、磁学和电学等途径\cite{alin2008}。机械振动、推进器运转和绕流可产生辐射噪声\cite{liu2019sail}，艇体及内部磁性材料的磁化可引起外部磁场变化\cite{bullard1961,bosser2014}。金属与海水之间的电化学过程则会形成水下电位分布\cite{thiel2020}。这些特征分别反映目标结构、运行状态及其与海水的相互作用。电学观测由此成为水下目标物理特征研究的一个分支。

\subsection{水下电学源的构成与探测困难}

舰船水下电学特征的研究源于对海洋腐蚀及其防护过程的认识。早期试验关注异种金属接触如何改变金属在海水中的腐蚀行为，随后的理论研究开始描述电偶作用在介质中形成的空间电位分布\cite{davy1824,waber1953}。在船体尺度上，金属与海水界面的腐蚀及阴极保护过程形成电流，周围水体呈现相应的电位分布\cite{thiel2020}。因此，水下电学特征的形成涉及金属界面、导电水体和电流回路之间的共同作用，其研究内容也从材料防护扩展到源的空间结构及其外部表现。

除相对稳定的电化学分量外，航行器运行及其与水体的相互作用还会产生时变的电学特征。船模试验观察到螺旋桨转动与阴极保护输出电流的低频调制有关\cite{akomah2002}。船载设备产生的杂散电磁场、外部杂散电流对船体电位的影响，以及尾流中的感应电场也分别得到研究\cite{yang2006,wang2017stray,chen2024wake}。这些现象具有不同的物理来源，源端电流变化、介质中的电位变化和接收器输出分别描述这一过程的不同环节。

理论计算和测量技术的发展，使水下电学研究逐步从源机制解释走向实际水体中的观测。已有研究以船体周围电位测量检验理论计算，也在受控水池中记录不同方向的电场分量\cite{zamani1988,jiang2024}。与此同时，水流感应电位、运动平台噪声和电极漂移等因素进入测量结果\cite{crona2001,liu2022motion,wang2025agcl}。测点位置与方向改变时，目标相关变化及背景的相对贡献也会改变。如何解释空间分布与衰减，如何区分源变化与接收响应变化，成为由模型走向实际观测时需要处理的问题。

\subsection{从测场到测电荷的观测量转变}

水下电学观测采用不同的接收原理，其直接记录的物理量也有所不同。电极式测量记录两点之间的电位差，并结合测量基线解释相应方向上的电场信息\cite{jiang2024}。微结构电学传感则提供了通过感应响应获取外部电学信息的另一类器件形式。已有 MEMS 微型静电场仪利用微机械结构调制感应电极的响应，建立器件输出与外部静电场之间的关系\cite{horenstein2001mems}。测量量、器件输出及其标定关系决定了观测结果的物理含义。

本文讨论的感知系统以敏感单元的电荷响应为感知对象，记录感应电荷对应的差分电压。源端关注电荷量、空间分布及其随时间的变化，介质中的电位与电场用于描述源与接收位置之间的空间关系，接收端则关注敏感单元上的电荷响应及电压输出。源电荷、接收电荷和差分电压是相互联系的不同物理量。本文所称极化带效应的物理定义为\underline{\hspace{6em}}，其与外部电学变化的对应关系为\underline{\hspace{6em}}，接收电荷与差分电压之间的标定关系为\underline{\hspace{6em}}。

沿这一观测关系，源的形成过程与接收器的响应过程需要分别描述。源端的反应和输运决定电学变化如何产生，介质条件与边界影响其空间分布，敏感结构及读出电路则决定最终记录的响应。源端反应与输运模型的依据及适用条件为\underline{\hspace{6em}}，接收端准静态描述及边界条件的依据为\underline{\hspace{6em}}。电荷源的表征与接收端的标定由此构成衔接物理解释和实际记录的两个环节。

\subsection{研究意义}

水下目标的电学特征联系着目标本身的物理过程及其与周围水体的相互作用，对这些特征的研究能够拓展水下观测所获取的信息。已有研究从电化学过程解释介质中的电位分布，并通过多通道测量记录不同位置和方向的电学变化\cite{thiel2020,jiang2024}。在此基础上，面向电荷响应开展感知研究，能够进一步联系源的形成、介质中的空间分布与敏感单元的输出，使观测结果与其物理来源建立更清晰的对应。本文围绕这一联系开展研究，其意义体现在水下电学特征的物理认识、感知响应的实验表征以及连续观测信息的利用等方面。

在理论认识方面，以电荷量、空间分布及其时间变化描述源特征，有助于将源的形成过程与接收端的观测结果联系起来。源端状态、介质场量和接收电荷响应分别承担不同的物理含义，明确三者之间的关系，可以为解释差分电压的变化提供依据，使幅值、时间尺度与空间差异对应到具体的物理问题。将源特征与观测响应纳入统一的描述，也有助于比较不同电学过程在接收端的表现，理解同类源在不同观测条件下为何呈现不同响应。这一研究能够深化对水下电学特征可观测性的认识，为电荷感知系统的响应解释和后续模型研究提供理论基础。

在实验方法方面，等效源为连接理论描述与实际测量提供了可控的研究对象。围绕具有代表性的源特征建立等效关系，可以将源的空间结构及时间变化落实为可重复的实验条件，使不同条件下的观测具有共同的比较依据。其价值在于将模型预测、源端记录和接收输出联系起来，通过重复测量考察这种联系是否稳定，并据此认识接收系统能够表征哪些变化。等效源与标定关系的结合，还能够为比较不同通道和不同实验批次的响应提供参照，推动电荷感知研究从现象记录发展到具有明确物理对象的实验表征。

在信号解释方面，接收记录包含源、环境和测量系统共同作用的结果，理解其中可重复变化的来源，是利用电学观测信息的重要环节。将物理模型与数据分析相结合，有助于把时域波动、频谱成分及多通道差异与实际观测条件对应起来，形成从信号变化到物理解释的联系。其中，模型为具有明确物理关系的响应提供解释，对照与重复记录则为复杂耦合现象的分析提供依据。这样的研究能够增强特征选择的物理依据，使信号处理所得的信息与源表征和实验检验相互衔接，并为评价特征的重复性与稳定性提供支撑。

在感知技术与应用方面，电荷响应研究为水下电学观测提供了新的信息获取视角。阵列中的时间变化与空间差异，可用于研究外部电学变化在不同敏感单元上的共同表现及其区别，为认识系统的感知能力提供实验基础。进一步将可重复特征与连续判别联系起来，可以把观测记录转化为对状态变化的判断，并通过检出、误报和延迟等指标评价信息利用的效果。这一研究为电荷感知技术由器件响应走向连续观测建立联系，也为海洋环境观测、水下装备状态感知及水下目标探测积累方法与测量依据。

\section{国内外研究现状}
\subsection{水下电学源机理与建模}

电学特征研究从海洋腐蚀试验逐步发展到介质电位的理论推导和船体电位测量。Davy 通过浸泡试验研究了海水中异种金属接触对铜护板腐蚀的影响\cite{davy1824}。Waber 将电偶腐蚀表述为电解质中的空间电位分布问题\cite{waber1953}，Zamani 随后将边界元法用于船体阴极保护的电位计算，并以原型船的测量数据对照计算结果\cite{zamani1988}。现代研究进一步用非线性极化模型或电位叠加方法描述腐蚀与防护电流形成的电位分布\cite{wang2018corrosion,thiel2020}。Schaefer 等基于潜没艇体模型研究了水下电位与水面以上电位特征之间的联系\cite{schaefer2019}。

除电化学直流分量外，已有研究还涉及轴系相关交流分量、设备杂散场和流动感应电场。Akomah 等通过船模试验研究了螺旋桨转动与阴极保护输出电流低频调制的关系\cite{akomah2002}。设备杂散场和外部杂散电流也有相应研究，前者来自船载设备，后者可干扰船体阴极保护电位\cite{yang2006,wang2017stray}。Chen 等在拖曳模型形成的尾流中测量了感应电场\cite{chen2024wake}。这些工作使用的观测量包括源端电流、介质电位和电场分量。分析电场来源时，需要把观测量与相应物理过程对应起来。

国内研究从腐蚀电学特征逐步拓展到运动与流动相关电学现象。哈尔滨工程大学的赵景波研究了舰船腐蚀电磁场的测量\cite{zhao2006}。海军工程大学的 Wang 等将非线性极化特性纳入边界元计算，研究船体腐蚀电场\cite{wang2018corrosion}。南京理工大学与南京航空航天大学的 Chen 等通过拖曳模型研究尾流中的感应电场\cite{chen2024wake}。这些工作分别从界面电化学条件和水体运动出发，扩展了对水下电学源的认识。

上述电位建模工作以介质电位及其空间变化为主要结果，并通过界面条件描述源的作用。若进一步讨论电荷如何产生和重新分布，还需考察界面反应与介质输运。以电荷分布为状态量的反应与输运研究及其代表性成果为\underline{\hspace{6em}}。从这些研究过渡到电荷响应观测，关键在于说明模型中的源变量、介质变量和接收变量怎样对应，并为这种对应提供可检验的依据。

\subsection{水下电场传感与测量技术}

水下电场的接触式传感研究主要关注电极对的输出噪声、响应灵敏度和长期电位稳定性。Wang 等对银／氯化银（Ag/AgCl）电极对分别开展了长期电位差、外加信号响应和噪声测量\cite{wang2025agcl}。Crona 等在同一潮汐水道比较了石墨纤维电极和 Ag/AgCl 电极。两类电极均记录到流动感应电位，所用石墨纤维样品对缓慢变化的响应有所衰减\cite{crona2001}。Zai 等比较改性与未改性碳纤维电极，在同一试验中观察到改性样品的电位差漂移和自噪声减小，对较弱的交流激励响应也更清楚\cite{zai2020carbon}。这些对照说明，电极材料、界面状态及信号频率都会影响测量结果。比较灵敏度、噪声和稳定性时需要采用一致的测量量、带宽及观测时长。

公开的 MEMS 电场传感研究提供了微结构读出电学响应的实例。Horenstein 和 Stone 研制了微型静电场仪，通过微机械结构调制感应电极的响应以读取外部静电场\cite{horenstein2001mems}。该研究的标定对象是静电场。本文所用电荷传感阵列的技术出处为\underline{\hspace{5em}}。该阵列与接触式电极在相同条件下的标定对照结果为\underline{\hspace{5em}}。

阵列测量利用不同测点或方向上的电响应表征空间分布。由多点观测进一步估计位置时，需把观测量与相应的场模型联系起来。受控水槽中的仿生电感知研究已利用电响应变化开展样品定位\cite{solberg2008}，相关结果为说明空间观测与位置估计的关系提供了实例。

多通道测量的发展使不同位置和方向的电学响应能够被同步记录。中国地质大学（北京）与武汉第二船舶设计研究所合作研制多通道水下电场测量系统，在水池中用受控移动信号源模拟船模电场，并记录三个方向的电场分量，验证系统的基本测量功能\cite{jiang2024}。与单个测点相比，多点记录提供了空间差异的信息，也将通道间响应的一致性带入测量问题。

综上，电极材料研究、微型器件研究与多通道系统研究分别改善或扩展了电学观测的不同环节。噪声、漂移和频率响应描述时间域内的测量性能，测点布置及通道响应则影响空间信息的获取。对于电荷响应阵列，还需明确其实际感知机制及输出的标定含义，才能将相关性能与既有电场测量结果进行对应。

\subsection{电荷源表征、反演与等效构造}

由源预测观测结果与由观测结果解释源，构成两个不同方向的问题。已有船体电位计算给出了源端条件到介质电位的预测实例\cite{zamani1988,wang2018corrosion}。反向解释则需要确定观测中哪些信息对应源的空间结构，哪些变化可能来自介质边界或接收器响应。关于电荷源反演的非唯一性、正则化及等效源表示的代表性研究为\underline{\hspace{6em}}。

受控源实验为模型和观测之间的比较提供了可重复的对象。Jiang 等在水池中以移动的受控信号源模拟船模电场，用已知实验布置检验测量系统\cite{jiang2024}。这一工作体现了通过构造源来研究接收响应的实验思路。以电荷量及其空间分布为等效对象的构造方法、适用条件和验证证据为\underline{\hspace{6em}}。

等效关系需要同时说明所保留的源特征及用于比较的接收量。介质电位相近、电荷分布相近和接收器输出相近分别对应不同层次的比较。对电荷响应感知而言，需要建立能够连接这些层次的物理关系，并通过重复观测检验这种联系。其研究重点是等效对象和验证量之间的一致性。

\subsection{电学信号处理与特征提取}

数字频谱分析的发展包含两个相互关联的问题，即如何高效计算离散频谱，以及如何利用有限长度的记录估计功率在频率上的分布。1965 年，Cooley 和 Tukey 通过分解变换长度，将离散傅里叶变换化为若干较短的变换，降低了频谱计算量\cite{cooley1965}。1967 年，Welch 用快速傅里叶变换计算各段记录的加窗周期图，再对各段估计取平均，以减小功率谱估计的波动\cite{welch1967}。离散傅里叶变换给出记录的频率成分，功率谱估计进一步描述功率随频率的分布，为分析周期信号与背景噪声提供了依据。

当频率成分随时间变化时，研究问题进一步涉及这些成分出现的时刻及其持续过程。1984 年，Grossmann 和 Morlet 利用同一波形的平移与伸缩构造小波族，给出了连续小波变换的数学表述。平移对应观测位置，伸缩对应分析尺度，因而可以描述信号在不同时间和尺度上的局部变化\cite{grossmannMorlet1984}。1989 年，Mallat 从多分辨率分析出发，将相邻分辨率之间的细节表示为正交小波系数，并给出了逐级计算方法\cite{mallat1989}。在海洋电场研究中，连续小波变换也用于对照去噪前后的实测时频图，观察轴频线谱在时间上的延续和保留情况\cite{wang2024vmd}。

信号分解研究进一步关注如何根据观测记录估计其中的成分。1998 年，Huang 等提出经验模态分解，按数据的局部时间尺度提取本征模态函数，以分析非平稳信号\cite{huang1998}。2014 年，Dragomiretskiy 和 Zosso 提出变分模态分解，将各模态及其中心频率一并纳入估计，并用这些模态共同重构输入信号\cite{dragomiretskiyZosso2014}。两种分解分别依据局部时间尺度和模态的频谱集中性，为后续研究在不同成分上处理噪声提供了途径。

海洋电场去噪研究进一步把模态分解与噪声抑制结合。2022 年，Liu 等先考察运动平台电场记录的噪声特征，再用改进的自适应噪声完备集合经验模态分解和自适应阈值处理进行去噪与重构\cite{liu2022motion}。2024 年，Wang 等将变分模态分解与多窗谱减结合，对模态中的噪声进行抑制，并用模拟信号和实测电场记录检验去噪效果及线谱保留情况\cite{wang2024vmd}。

从谱线表示到信号判别，还需要将频谱特征转化为可比较的判据。1998 年，Chen、Donoho 和 Saunders 提出的基追踪方法，以系数的 $\ell_1$ 范数最小为准则，选择能够表示原信号的字典分解\cite{chen1998basis}。2025 年，Hu 等将基追踪引入低频分析与记录方法（low-frequency analysis and recording，LOFAR），估计实测电场记录的低频谱，并在同一海洋电场记录上比较不同频谱估计方法所得的时频表示与阈值判别结果\cite{hu2025lofar}。

这些方法分别提供频率组成、局部时频变化、模态成分和稀疏表示的信息。现有电场应用研究展示了运动噪声处理、轴频线谱保留及低频谱估计的具体结果。电荷响应记录中的有效特征还需要与其实际信号形态对应，并以原始记录和已知实验条件检验处理后的变化。谱线、宽带波动和非平稳变化涉及不同的表征问题，方法的选择应与所分析的信号成分相联系。

\subsection{水下目标检测与识别}

电场检测判据要明确检出与误报的统计含义。Neyman 和 Pearson 的假设检验理论以控制第一类错误概率为前提，比较对备择假设的检验功效\cite{neymanPearson1933}。恒虚警检测（constant false alarm rate，CFAR）研究进一步关注背景水平变化时虚警水平的控制\cite{finnJohnson1968}。在电场研究中，检出率应以已知目标状态的记录为依据，误报率则对应仅有背景的记录。由谱线增强所得的时频图与由判据所得的检出结果分别回答信号能否辨认和判决是否可靠的问题。

电学观测的识别研究还涉及从信号中估计样品的属性。Bai 等在受控水槽的主动电感知实验中比较不同样品的幅值与相位响应，观察到部分幅值响应接近的样品可由相位信息区分\cite{bai2012}。后续研究用监督回归估计水槽样品的几何属性\cite{bai2015}。采用被动电场实测数据开展舰船类别识别的代表性研究为\underline{\hspace{5em}}。评价相关识别方法时，需说明信号来源、类别标签及测试记录的独立性。

连续观测将判别对象从一个记录片段扩展到随时间变化的状态。相应研究需要说明状态变化的时刻、判断的持续性及错误判决的分布。变化点检测与在线状态判别的代表性研究为\underline{\hspace{6em}}。从谱线表示、样品属性估计到连续状态判别，已有方法的输入、输出和验证任务各不相同。建立检测评价时，明确观测对象和参考状态有助于比较不同方法的实际作用。

\subsection{研究现状总结与存在问题}

水下电学研究已经形成源机理、介质电位计算、传感测量和信号分析等相互联系的研究方向。对于以电荷响应为感知对象的系统，首先需要明确源端状态量与接收端观测量的对应关系。现有电位计算能够解释特定边界条件下的空间分布，电荷源的表征还需要进一步说明这些条件如何联系源的形成与变化。

其次，理论描述与等效实验之间需要具有一致的比较对象。受控源能够提供可重复的激励，但从源特征到接收输出的对应需要结合介质、边界和器件响应解释。等效构造与接收标定应共同支持这种对应，而非分别停留在各自的物理量上。

再次，实际观测中包含设备运行、水体运动及背景变化的共同影响。源相关特征的解释需要结合重复记录和对照条件，将信号的统计变化与可能的物理来源相联系。最后，连续判别需要把特征提取与有参考状态的评价衔接起来，使检出、误报和延迟具有明确的统计对象。这些问题形成本文开展源表征、实验观测与判别研究的出发点。

\section{关键科学问题与研究内容}
\subsection{关键科学问题}

本文围绕电荷源、观测映射和响应判别三个相互衔接的问题展开。源端问题关注具有代表性的水下电学过程能够用哪些电荷量、空间结构及时间变化来表征，以及这些特征如何在可控条件下复现。源的选择依据物理代表性、观测关联性与重复构造的可行性，阴极保护作为电化学源的一个研究实例。

观测映射问题关注源特征如何在介质和边界作用下表现为接收端的电荷响应，并进一步形成差分电压记录。研究需要区分源参数变化、空间条件变化和接收响应变化的作用，明确理论量与实验量之间的联系。源电荷与阵列输出之间的具体物理关系为\underline{\hspace{6em}}。

响应判别问题关注背景漂移、设备干扰及流动相关变化共存时，源相关变化能否在连续观测中被稳定辨认，以及其解释能否得到对照实验支持。探测任务由存在性判断进一步延伸到方位和距离信息的获取，各层级分别对应不同的观测要求。本文关于方位与距离估计的研究范围为\underline{\hspace{6em}}。

\subsection{研究内容}

围绕电荷源的表征与等效构造，梳理代表性电学过程及其可观测特征，明确电荷量、空间分布与时间变化的描述方式。通过理论分析联系源端条件与实验中的可控量，为构造具有可重复特征的等效源提供依据。

围绕源、介质与接收之间的关系，区分源端形成过程、介质中的空间分布和器件读出过程。理论计算与仿真用于描述具有明确物理条件的部分，实验观测用于检验预测与响应之间的联系。具体接收模型及标定参数为\underline{\hspace{6em}}。

围绕实验表征，按照激励可控、条件可记录、测量可重复和背景可对照的原则组织水槽观测。记录源端条件与接收端输出，比较重复测量中的共同特征及差异，将实验结果与等效源的理论描述相联系。

围绕特征提取与连续判别，采用时域、频域和时频分析描述响应变化，并结合多通道信息研究其时间一致性和空间差异。能够由明确物理关系解释的成分采用模型分析，其余耦合响应结合对照数据进行归因，两者在同一等效源及观测条件下比较。以连续记录为输入形成状态判别，并用检出、误报和延迟评价其表现。

\subsection{技术路线}

研究从电学源的物理表征出发，建立源、介质与接收响应之间的联系，再以等效源连接模型与实验。实验记录用于检验响应关系并提取可重复特征，连续判别用于评价这些特征对状态变化的指示能力。模型解释和数据分析共同服务于观测量的物理归属及判别结果的验证。

\section{论文结构安排}
% 按用户要求保留标题，待各章完成后补写。

\section{本章小结}

本章从海洋观测需求和多物理场探测的发展引出水下电学研究，区分了源电荷、介质场量及接收电荷响应的物理含义。通过梳理源机理、传感测量、等效构造、信号处理与检测识别的研究，提出源表征、观测映射和响应判别三个问题，并说明其与理论分析、可控实验及连续观测评价的联系。后续章节围绕典型电荷源及其等效表征展开。
